\section{Lecture 10: Analytic (Tate) Adic Spaces}
\label{lec:10}

This lecture turns to the comparison with Huber's
classical theory of \emph{adic spaces}, in the special ``analytic'' --- here
called \emph{Tate} --- case. There are two halves. First, we recall Tate Huber
rings (Banach algebras with a topologically nilpotent unit), the notion of a
\emph{sheafy} Huber pair, and the theorem of Kedlaya--Liu that finite projective
modules glue to vector bundles when $(A,A^{+})$ is sheafy. Second, and this is
the point of the lecture, we re-prove and vastly generalise that theorem inside
the solid formalism of \cref{lec:7,lec:9}: for \emph{any} Huber pair with $A$
Tate --- with no sheafiness hypothesis --- pseudocoherent complexes, perfect
complexes and vector bundles all glue over the adic spectrum, because they are
detected inside the (always well-behaved) sheaf of solid derived categories. The
argument reduces the classical analysis to two abstract finiteness conditions,
\emph{pseudocoherence} and \emph{nuclearity}, whose intersection is pinned down
by a Fredholm-theoretic argument that is the one genuinely analytic step.

Throughout, ``ring'' means commutative ring, everything lives in light condensed
abelian groups, and we use the solid theory $\Solid\subseteq\CondAbl$
(\cref{lec:5,lec:6}), the relative theory over $\mathbb{Z}[T]$ (\cref{lec:7}),
the abstract notion of analytic ring (\cref{lec:8}) and the localization of
solid analytic rings over the valuative/adic spectrum (\cref{lec:9}).

\subsection{Tate adic spaces}

\begin{remark}[A conflict of terminology]
\label{rmk:10:terminology}
Every adic space gives rise to an analytic space in our sense (\cref{lec:9}).
But within the theory of adic spaces there is already an established qualifier
``analytic'' for a certain subclass. To avoid overloading the word, in this
lecture Scholze calls the members of that subclass \emph{Tate} adic spaces.
\end{remark}

\begin{definition}[Tate Huber ring; Tate adic space]
\label{def:10:tate}
A Huber ring $A$ (\cref{def:8:huber-ring}) is \emph{Tate} if it has a
topologically nilpotent unit $\pi\in A$, a so-called \emph{pseudo-uniformizer}.
An adic space $X$ is \emph{Tate} (the property classically called
\emph{analytic}) if it is covered by affinoids $\Spa(A,A^{+})$ with $A$ Tate.
Equivalently,
\begin{equation*}
X\ \text{is Tate}\quad\Longleftrightarrow\quad
\text{the completed residue fields at all points of }X\ \text{are not discrete}
\end{equation*}
(so that they are honest nonarchimedean fields, rather than discrete rings).
\end{definition}

\begin{example}[Tate rings]
\label{ex:10:tate-examples}
$\mathbb{Q}_p$ is Tate, with pseudo-uniformizer $\pi=p$; more generally any
nonarchimedean local field, and any Huber ring over such a field, is Tate ---
once a topological unit is present downstairs, it remains a topological unit
after any base change up.
\end{example}

\begin{remark}[Intuition: schemes, formal schemes, generic fibres]
\label{rmk:10:intuition}
The heuristic is that adic spaces interpolate between schemes and formal schemes.
A scheme sits inside a formal scheme; the associated adic space contains, at one
extreme, the ``scheme-like'' locus, and at the other extreme the Tate locus,
which is morally the generic fibre of a formal scheme with all scheme-like
points removed. The basic picture is the formal scheme $\Spf\mathbb{Z}_p$, i.e.
$\Spa(\mathbb{Z}_p,\mathbb{Z}_p)$, with its closed (scheme-like) point
$\Spec\mathbb{F}_p$ and its open Tate generic fibre $\Spa(\mathbb{Q}_p,\mathbb{Z}_p)$:
% board 3 @ 00:08:00--00:09:06
\[
\begin{tikzcd}[column sep=large]
\Spec \mathbb{F}_p \arrow[r, hook] & \Spa(\mathbb{Z}_p,\mathbb{Z}_p) & \Spa(\mathbb{Q}_p,\mathbb{Z}_p) \arrow[l, hook']
\end{tikzcd}
\]
The closed point is scheme-like; the open complement is Tate.
\end{remark}

\subsection{Structure of Tate Huber rings}

\begin{proposition}[Tate Huber rings are Banach algebras]
\label{prop:10:banach-structure}
Let $A$ be a Tate Huber ring, $\pi\in A$ a pseudo-uniformizer, and
$A_{0}\subseteq A$ a ring of definition. One may always arrange $\pi\in A_{0}$
(and that $A_{0}$ contains any prescribed power-bounded element); then $A_{0}$
carries the $\pi$-adic topology --- a priori it could carry the $I$-adic topology
for some finitely generated ideal $I$, but for a Tate ring the topology is always
$\pi$-adic. This turns $A$ into a Banach algebra with unit ball $A_{0}$, via
\begin{equation}
\label{eq:10:norm}
|f|\ =\ 2^{-n},\qquad n\ =\ \sup\{\,m\in\mathbb{Z}\ \mid\ \pi^{-m}f\in A_{0}\,\},
\end{equation}
so that $|\pi|=\tfrac12$ and $|\pi^{-1}|=2$.
\end{proposition}

\begin{remark}[Equivalence with ultrametric Banach rings]
\label{rmk:10:banach-equivalence}
The base $2$ and the choice of $\pi$ in \eqref{eq:10:norm} are arbitrary; only the
induced topology matters. Passing to that topology gives an equivalence of
categories
\begin{equation*}
\{\text{Tate Huber rings}\}\ \simeq\
\left\{\begin{array}{c}
\text{complete ultrametric Banach rings } B \text{ admitting } \pi\in B\\
\text{with } 0<|\pi|<1 \text{ and } |\pi|\,|\pi^{-1}|=1
\end{array}\right\},
\end{equation*}
with continuous maps as morphisms. On the right one should not fix a norm but only
require that \emph{some} equivalent norm has the displayed multiplicativity on
$\pi$. (The zero ring is a degenerate exception to the naive statements about
pseudo-uniformizers, and is best set aside.)
\end{remark}

\subsection{Sheafiness and the theorem of Kedlaya--Liu}

Recall (\cref{lec:9}) that Huber equips $\Spa(A,A^{+})$ with a structure
presheaf $\Ohub$; its defect is that it is not always a sheaf.

\begin{definition}[Sheafy]
\label{def:10:sheafy}
A Huber pair $(A,A^{+})$ is \emph{sheafy} if the Huber structure presheaf
$\Ohub_{\Spa(A,A^{+})}$ is a sheaf.
\end{definition}

\begin{theorem}[Kedlaya--Liu \cite{kedlaya2015relativepadichodgetheory}]
\label{thm:10:kedlaya-liu}
Let $A$ be Tate and $(A,A^{+})$ sheafy. Then base change to the structure sheaf
is an equivalence
\begin{equation}
\label{eq:10:kedlaya-liu}
\FProj(A)\ \xrightarrow{\ \sim\ }\ \VB\bigl(\Spa(A,A^{+}),\,\Ohub_{\Spa(A,A^{+})}\bigr),
\qquad M\longmapsto M\otimes_{A}\Ohub_{\Spa(A,A^{+})},
\end{equation}
between finite projective $A$-modules and vector bundles (locally free sheaves of
finite rank) on the adic spectrum. In other words, finite projective modules
glue. There is a version for (pseudo)coherent modules, whose precise statement is
more subtle.
\end{theorem}

\begin{remark}[History; the plan for today]
\label{rmk:10:history}
Scholze found this theorem striking when it appeared: one expected that in
situations where sheafiness \emph{could} be proved good things would follow, but
not that sheafiness alone --- the property most prone to fail --- could be the
sole hypothesis of a nontrivial theorem. The original Kedlaya--Liu proof is a
delicate change-of-basis argument that ``does some actual analysis''. Kremnizer
had suggested to Scholze already at the time that, working in derived categories
with the right notion of modules, one can simply glue; the task today is to make
that concrete and recover the vector-bundle statement. The reference for the
solid treatment is Andreychev
\cite{andreychev2021pseudocoherentperfectcomplexesvector}; the same argument
proves the analogue of the corresponding statement in the complex-analytic
setting of \cite{clausen2020analytic}. The lecture has two parts:
\begin{enumerate}
\item the relation between Huber's $\Ohub$ and the solid structure sheaf $\mathcal O$
of \cref{lec:9}, culminating in the fact that \emph{when} $(A,A^{+})$ is sheafy
the two agree;
\item the gluing of pseudocoherent and perfect complexes, hence of vector
bundles, for an \emph{arbitrary} Tate Huber pair.
\end{enumerate}
\end{remark}

\subsection{Huber's presheaf versus the solid structure sheaf}

Fix a Huber pair $(A,A^{+})$ with $A$ Tate. Throughout we use the adic spectrum
$\Spa(A,A^{+})$ of \emph{continuous} valuations, so that in forming a rational
subset $U(f_{1},\dots,f_{n}/g)$ we assume $(f_{1},\dots,f_{n})$ generates an open
ideal --- which, $A$ being Tate, is automatically the unit ideal, so
$(f_{1},\dots,f_{n})=A$.

\begin{proposition}[Huber's structure presheaf]
\label{prop:10:huber-presheaf}
For a rational subset $U=U(f_{1},\dots,f_{n}/g)\subseteq\Spa(A,A^{+})$ (so
$(f_{1},\dots,f_{n})=A$),
\begin{equation}
\label{eq:10:huber-presheaf}
\Ohub(U)\ =\ A\langle T_{1},\dots,T_{n}\rangle\big/\,\overline{(gT_{1}-f_{1},\dots,gT_{n}-f_{n})},
\end{equation}
the convergent power series algebra in $T_{1},\dots,T_{n}$ modulo the
\emph{closure} of the ideal $(gT_{i}-f_{i})$. Here $g$ is already invertible
before quotienting (its ideal contains $f_{1},\dots,f_{n}$, hence is the unit
ideal), and $T_{i}$ is the new function $f_{i}/g$, bounded by $1$ on $U$.
\end{proposition}

The point of Huber's construction is that a quotient of a Banach algebra by a
\emph{closed} ideal is again a Banach algebra --- so one takes the closure and
stays inside the world of Huber rings. The solid structure sheaf of \cref{lec:9}
instead uses a \emph{derived} quotient, and this is where the comparison lives.

\begin{notation}[Derived quotient/Koszul complex]
\label{not:10:koszul}
For a condensed ring $R$ and $x_{1},\dots,x_{n}\in R$, write
\begin{equation*}
R/^{\mathbb{L}}(x_{1},\dots,x_{n})\ =\ R\otimes^{\mathbb{L}}_{\mathbb{Z}[x_{1},\dots,x_{n}]}\mathbb{Z}
\end{equation*}
for the derived quotient, an animated condensed ring, computed by the Koszul
complex
\begin{equation*}
\Bigl(\,\textstyle\bigwedge^{n}(\textstyle\bigoplus_{i=1}^{n}R)\to\cdots\to
\bigwedge^{2}(\textstyle\bigoplus_{i=1}^{n}R)\to
\textstyle\bigoplus_{i=1}^{n}R\xrightarrow{\,(x_{i})_{i}\,}R\,\Bigr).
\end{equation*}
\end{notation}

\begin{remark}[The solid value of a rational open]
\label{rmk:10:solid-value}
The solid structure sheaf assigns to $U$ the endomorphisms of the unit of the
analytic ring $(\mathcal O(U),\mathcal O^{+}(U))^{\solid}$; concretely
$\mathcal O(U)$ is $A[\tfrac1g]$ solidified so that each $f_{i}/g$ is
power-bounded. Its value has the same shape as \eqref{eq:10:huber-presheaf}, but
with a derived quotient in place of Huber's closure. In two steps
(cf.\ \cref{lec:7}):
\begin{itemize}
\item start with the derived localization
$A[\tfrac1g]=A[T_{1},\dots,T_{n}]/^{\mathbb{L}}(gT_{i}-f_{i})$ (the sequence
$gT_{i}-f_{i}$ is regular, so ``derived'' does no harm), and then
$(T_{1},\dots,T_{n})$-solidify. Solidifying the polynomial variables produces the
Tate algebra,
\begin{equation*}
A[T_{1},\dots,T_{n}]^{\mathbb{L}(T_{1},\dots,T_{n})\text{-}\solid}\ =\
A\langle T_{1},\dots,T_{n}\rangle,
\end{equation*}
because joining a variable and solidifying commute with the presentation
$A=A_{0}[\tfrac1\pi]$, $A_{0}=\varprojlim_{n}A_{0}/\pi^{n}$, as limits and
colimits (\cref{lec:7}).
\end{itemize}
\end{remark}

\begin{corollary}[Huber $=$ quasiseparation of $\pi_{0}$ of the solid value]
\label{cor:10:huber-quasisep}
With $\mathcal O(U)$ the solid value as above,
\begin{equation*}
\Ohub(U)\ =\ \bigl(\pi_{0}\mathcal O(U)\bigr)^{\mathrm{qs}},
\end{equation*}
Huber's ring is the \emph{quasiseparation} of the underlying static condensed ring
$\pi_{0}\mathcal O(U)=A\langle T_{1},\dots,T_{n}\rangle/(gT_{i}-f_{i})$.
\end{corollary}

\begin{remark}[The quasiseparation functor]
\label{rmk:10:qs-functor}
The inclusion of quasiseparated condensed sets into all condensed sets,
$\mathrm{CondSet}^{\mathrm{qs},\mathrm{light}}\subseteq\CondSetl$, has a left adjoint
\begin{equation*}
X\ \longmapsto\ X^{\mathrm{qs}},
\end{equation*}
a ``maximal Hausdorff quotient''. For non-obvious reasons it commutes with finite
products, hence preserves algebraic structures (being a group, ring, etc.), so
$(\pi_{0}\mathcal O(U))^{\mathrm{qs}}$ is again a ring; on a light condensed set
the quasiseparation is again light. Passing from $\pi_{0}\mathcal O(U)$ to its
quasiseparation is exactly Huber's passage to the closure of the ideal
$(gT_{i}-f_{i})$.
\end{remark}

\begin{lemma}[Quasiseparation via equivalence relations]
\label{lem:10:qs-equiv}
Write a condensed set $X$ as a quotient $X=\widetilde X/R$ of a quasiseparated
$\widetilde X$ (e.g. a profinite set surjecting onto it) by an equivalence
relation $R\subseteq\widetilde X\times\widetilde X$. Then there is a minimal
quasicompact injection $\overline R\subseteq\widetilde X\times\widetilde X$ that
is again an equivalence relation and contains $R$, and
\begin{equation*}
X^{\mathrm{qs}}\ =\ \widetilde X/\overline R .
\end{equation*}
One obtains $\overline R$ by intersecting all quasicompact-injective equivalence
relations containing $R$ (a transfinite closure that may add more than the naive
closure of $R$). If $X$ is a group one may take $\widetilde X$ a group and $R$ a
subgroup, in which case $\overline R$ is the closure of the subgroup.
\end{lemma}

\begin{lemma}[Closed subsets are quasicompact injections]
\label{lem:10:closed-subsets}
Let $X$ be any sequential topological space (so that $\underline X$ is a light
condensed set, \cref{lec:3}). Then
\begin{equation*}
\{\text{closed subsets of }X\}\ \xrightarrow{\ \sim\ }\
\{\text{quasicompact injections into }\underline X\}.
\end{equation*}
Indeed the pullback of a closed subset $Z\subseteq X$ along any map from a
profinite set is a closed subset of a profinite set, hence profinite, hence
quasicompact; the quotient topology on $X$ makes this reversible.
\end{lemma}

We can now state the sharp comparison. It is essentially due to Kedlaya, phrased
here in the present language.

\begin{theorem}[Sheafy $\Leftrightarrow$ Huber $=$ solid]
\label{thm:10:sheafy-iff}
For a Tate Huber pair $(A,A^{+})$,
\begin{equation*}
\begin{aligned}
(A,A^{+})\ \text{sheafy}&\quad\Longleftrightarrow\quad \Ohub=\mathcal O\\
&\quad\Longleftrightarrow\quad
\forall\,U\ \text{rational},\ \mathcal O(U)\ \text{sits in degree }0
\ \text{and is quasiseparated}.
\end{aligned}
\end{equation*}
That is, whenever Huber's presheaf happens to be a sheaf, it agrees with the solid
structure sheaf $\mathcal O$ --- so the good properties of $\mathcal O$ transfer,
which is why sheafiness alone forces good behaviour.
\end{theorem}

\begin{proof}[Sketch]
The implication ``$\Leftarrow$'' is clear: if $\mathcal O(U)$ sits in degree $0$
and is quasiseparated for every rational $U$, then by
\cref{cor:10:huber-quasisep} it equals $\Ohub(U)$, which is therefore the (sheafy)
solid structure sheaf. For ``$\Rightarrow$'' it suffices to check the good
properties of $\mathcal O(U)$ after possibly refining the cover. One always
refines to \emph{simple Laurent covers} $\{|f|\le1\}\cup\{|f|\ge1\}$, and the
problem reduces to \cref{lem:10:key} below.
\end{proof}

\begin{lemma}[Key Lemma]
\label{lem:10:key}
Let $A$ be a Tate Huber ring and $f\in A$. Assume the Laurent-cover sequence
\begin{equation}
\label{eq:10:laurent-sequence}
0\ \to\ A\ \to\ A\langle \tfrac f1\rangle\oplus A\langle\tfrac1f\rangle\ \to\
A\langle \tfrac f1,\tfrac1f\rangle\ \to\ 0
\end{equation}
is exact (as it is, e.g., when $A$ is sheafy) --- here
$A\langle\tfrac f1\rangle=\Ohub(\{|f|\le1\})$,
$A\langle\tfrac1f\rangle=\Ohub(\{|f|\ge1\})$ and
$A\langle\tfrac f1,\tfrac1f\rangle=\Ohub(\{|f|=1\})$. Then Huber's rings agree
with the derived quotients,
\begin{equation*}
A\langle \tfrac f1\rangle=A\langle T\rangle/^{\mathbb{L}}(T-f),\qquad
A\langle \tfrac1f\rangle=A\langle T\rangle/^{\mathbb{L}}(1-Tf),
\end{equation*}
i.e.\ multiplication by $T-f$ (resp.\ $1-Tf$) on $A\langle T\rangle$ is injective
with closed image.
\end{lemma}

\begin{proof}[Idea]
One must see that $A\langle T\rangle\xrightarrow{\,T-f\,}A\langle T\rangle$ (resp.\
$1-Tf$) is injective with closed image. Using the exactness
\eqref{eq:10:laurent-sequence} the question may be analyzed over the two
localizations $A\langle\tfrac f1\rangle$ and $A\langle\tfrac1f\rangle$ --- in
other words one may assume $|f|\le1$ or $|f|\ge1$. In the case $|f|\le1$ the ideal
$(T-f)$ is closed because one successively kills the highest coefficient; the case
$|f|\ge1$ amounts to inverting $f$ already at the level of $A_{0}$, after which the
other localization becomes a tautology. Either way the statement can be checked by
hand. This is the funny reduction, due to Kedlaya, by which sheafiness of a single
Laurent cover controls everything.
\end{proof}

\subsection{Gluing pseudocoherent and perfect complexes}

We now drop the sheafiness hypothesis entirely. The classical finiteness notions
are those of \cite{Berthelot_1971} (SGA~6).

\begin{definition}[Pseudocoherent, perfect; SGA~6 \cite{Berthelot_1971}]
\label{def:10:pseudocoherent}
Let $R$ be any ring. A complex $C\in D(R)$ is \emph{pseudocoherent} if it can be
represented by a complex
\begin{equation*}
\cdots\to C_{n}\to\cdots\to C_{1}\to C_{0}\to\cdots\to C_{-m}\to 0\to 0\to\cdots
\end{equation*}
that is bounded to the right (zero in low degrees) but possibly unbounded to the
left, with each $C_{i}$ a finite free (equivalently, finite projective)
$R$-module. It is \emph{perfect} if such a representative can be chosen bounded on
both sides. Over a noetherian ring pseudocoherence says each homology group is
coherent; over a general ring it captures ``finitely generated with successively
finitely many relations''.
\end{definition}

\begin{theorem}[Gluing of finiteness conditions]
\label{thm:10:gluing}
Let $(A,A^{+})$ be any Huber pair with $A$ Tate. Then each of the assignments
sending a rational $U\subseteq\Spa(A,A^{+})$ to
\begin{equation}
\label{eq:10:gluing}
D^{\mathrm{pcoh}}(\mathcal O(U)(\ast)),\quad
D^{\mathrm{pcoh}}_{\ge n}(\mathcal O(U)(\ast)),\quad
\Perf(\mathcal O(U)(\ast)),\quad
\Perf^{[a,b]}(\mathcal O(U)(\ast))
\end{equation}
--- pseudocoherent complexes, pseudocoherent complexes in degrees $\ge n$,
perfect complexes, and perfect complexes admitting a representative in degrees
$[a,b]$, over the underlying static ring $\mathcal O(U)(\ast)$ --- is a sheaf of
$\infty$-categories. The special case $\Perf^{[0,0]}=\VB$ recovers the
Kedlaya--Liu theorem \eqref{eq:10:kedlaya-liu}, now without any sheafiness
hypothesis.
\end{theorem}

\begin{remark}[Warning: discreteness does not globalize]
\label{rmk:10:warning}
It is essential to work with the finiteness \emph{subcategories} of the solid
derived category rather than with the condition of being a discrete module. The
condition $M\in D(A(\ast))\subseteq D((A,A^{+})_{D})$ --- being (a complex of)
discrete $A(\ast)$-modules --- does \emph{not} globalize.
% board 18--19 @ 01:17:15--01:22:16
Consider the Tate elliptic curve, the quotient of the adic $\mathbb{G}_{m}$ over
$\mathbb{Q}_{p}$ by multiplication by $p$,
\begin{equation*}
a\colon\ \mathbb{G}_{m,\mathbb{Q}_p}^{\mathrm{ad}}\ \longrightarrow\
\mathbb{G}_{m,\mathbb{Q}_p}^{\mathrm{ad}}/p^{\mathbb{Z}},
\end{equation*}
the nonarchimedean analogue of $\mathbb{C}^{*}/q^{\mathbb{Z}}$. Let
$\Spa(A,A^{+})$ be a large open subset of the quotient --- remove a small disk
around one point, but keep it large enough to carry a global monodromy, and so
that it genuinely lies in the quotient rather than lifting to
$\mathbb{G}_{m}$. Then the lower-shriek $a_{!}\mathcal O^{\mathrm{ad}}_{\mathbb{G}_{m,\mathbb{Q}_p}}$,
restricted to $\Spa(A,A^{+})$, defines an object of $D((A,A^{+})_{D})$ that is
\emph{locally} discrete --- locally the covering map $a$ splits, so
$a_{!}\mathcal O$ is a direct sum of copies of the base indexed by
$\mathbb{Z}$ --- but is \emph{not} globally discrete: it does not lie in
$D(A(\ast))$. The monodromy obstructs the global splitting.
\end{remark}

\begin{remark}[Outline of the argument]
\label{rmk:10:outline}
By \cref{lec:9} we already know the assignment
\begin{equation*}
U\ \longmapsto\ D\bigl((\mathcal O(U),\mathcal O^{+}(U))_{D}\bigr)
\end{equation*}
is a sheaf of $\infty$-categories, and all the finiteness categories of
\eqref{eq:10:gluing} are cut out by finiteness conditions inside it. So gluing
reduces to a \emph{membership} statement: if $M\in D((A,A^{+})_{D})$ is such that,
over a cover, each base change
$M\otimes^{\mathbb{L}}_{(A,A^{+})_{D}}(\mathcal O(U),\mathcal O^{+}(U))_{D}$ lies
in one of the subcategories, then so does $M$ globally. This is carried out in
three steps: a pseudocoherent subcategory that globalizes formally, a
\emph{nuclear} subcategory that globalizes because internal $\Homop$ commutes with
the localizations, and the identification of the discrete pseudocoherent complexes
as the intersection of the two.
\end{remark}

\subsubsection*{Step 1: pseudocoherent complexes globalize}

\begin{definition}[Solid pseudocoherent complexes]
\label{def:10:solid-pcoh}
Define $D^{\mathrm{pcoh}}((A,A^{+})_{D})\subseteq D((A,A^{+})_{D})$ to consist of
the complexes representable by
\begin{equation*}
\cdots\to C_{n}\to\cdots\to C_{1}\to C_{0}\to\cdots\to C_{-m}\to 0,\qquad
\text{all } C_{i}=P,
\end{equation*}
bounded to the right, where $P$ is the compact projective generator (the free
complete module on a null sequence), subject to the finiteness condition that for
all $i\in\mathbb{Z}$ the functor $\Ext^{i}(C_{i},-)\colon\Mod(A)_{D}\to\mathrm{Ab}$
commutes with filtered colimits.
\end{definition}

That being in this subcategory globalizes is formal: it is a right-boundedness
statement together with the finite-Tor-dimension of the localizations, so it
descends.

\subsubsection*{Step 2: nuclear modules}

The eigenvector for detecting vector bundles requires a second class.

\begin{definition}[Nuclear modules]
\label{def:10:nuclear}
Let $P$ be the compact projective generator. An object
$C\in D((A,A^{+})_{D})$ is \emph{nuclear} if the natural map
\begin{equation}
\label{eq:10:nuclear}
C\otimes^{\mathbb{L}}_{(A,A^{+})_{D}}\underline{\Homop}(P,A)\
\xrightarrow{\ \sim\ }\ \underline{\Homop}(P,C)
\end{equation}
is an isomorphism --- ``all maps from $P$ into $C$ are trace-class''. Write
$D^{\mathrm{nuc}}((A,A^{+})_{D})$ for the full subcategory of nuclear objects.
\end{definition}

\begin{remark}[Provenance and generation]
\label{rmk:10:nuclear-remarks}
The name is borrowed, in spirit rather than to the letter, from Grothendieck's
nuclear spaces in functional analysis. Nuclear modules are generated under
colimits by the module $c_{0}(\mathbb{N},A)=A\langle T\rangle$ of $A$-valued null
sequences (and its shifts); in general by sequential colimits
$\varinjlim(Q_{0}\to Q_{1}\to\cdots)$ in which each transition is trace-class,
i.e.\ arises from an element of $Q_{i}^{\vee}\otimes Q_{i+1}$. Over $\mathbb{Q}_p$
they are generated by one-dimensional Banach spaces, still a large class. Nuclear
modules form a dualizable category --- a universal dualizable approximation to the
whole category --- which is one way to define analytic $K$-theory spaces
(Scholze noted this is not over a field).
\end{remark}

\begin{remark}[Why nuclearity globalizes]
\label{rmk:10:nuclear-globalize}
Both sides of \eqref{eq:10:nuclear} localize: the right by the same argument
Clausen used in \cref{lec:9} --- the different localizations commute because they
are all given by internal $\Homop$ out of fixed objects, and
$\underline{\Homop}(A,\underline{\Homop}(B,-))=\underline{\Homop}(A\otimes B,-)$.
The key point is that $\underline{\Homop}(P,-)$ commutes with the localizations of
interest; for a rational (rather than a general algebraic) localization this holds
because such a localization is itself given by an internal $\Homop$. Hence
nuclearity, being a local isomorphism condition, globalizes.
\end{remark}

\subsubsection*{Step 3: the intersection, and the Fredholm step}

\begin{proposition}[Discrete pseudocoherent $=$ solid pseudocoherent $\cap$ nuclear]
\label{prop:10:intersection}
Base change along $g\colon A(\ast)\to(A,A^{+})_{D}$ identifies
\begin{equation*}
D^{\mathrm{pcoh}}(A(\ast))\ \xrightarrow{\ \sim\ }\
D^{\mathrm{pcoh}}((A,A^{+})_{D})\ \cap\ D^{\mathrm{nuc}}((A,A^{+})_{D}).
\end{equation*}
\end{proposition}

\begin{proof}[Sketch]
Take $C$ in the intersection, represented as a pseudocoherent complex all of whose
terms are the generator $P$, and (by shifting) starting in degree $0$:
$\cdots\to P\to P\to0$. There is a tautological map $P\to C$, the identity in
degree $0$. Nuclearity makes this map trace-class, so it factors through an
approximation of $C$; concretely, the map $P\to C$ factors through the two-term
complex $[P\xrightarrow{\,1-g\,}P]$ for some trace-class endomorphism $g$ of $P$:
% board 23 @ 01:41:00--01:46:09
\[
\begin{tikzcd}[row sep=large, column sep=large]
& {[P\xrightarrow{\,1-g\,}P]} \arrow[d] \\
P \arrow[ur] \arrow[r] & C
\end{tikzcd}
\]
It then suffices to show $[P\xrightarrow{1-g}P]$ is perfect over $A(\ast)$. This is
classical Fredholm theory: with $g$ trace-class, $1-g$ is a Fredholm operator
(the identity is Fredholm, and a compact --- in particular trace-class ---
perturbation of a Fredholm operator is Fredholm), so its kernel and cokernel are
finite-dimensional and the kernel is a perfect complex. Iterating, one truncates
$C$ to a perfect complex of discrete modules. This is the single genuinely
analytic step of the whole argument.
\end{proof}

\begin{remark}[Conclusion]
\label{rmk:10:conclusion}
Combining the three steps: the solid pseudocoherent and nuclear subcategories both
glue, and their intersection is the discrete pseudocoherent (in particular
perfect, and $\Perf^{[0,0]}=\VB$) complexes, which therefore glue as well. The
remaining superscripted variants of \eqref{eq:10:gluing} glue by the same, easier,
bookkeeping. Thus vector bundles glue over $\Spa(A,A^{+})$ for every Tate Huber
pair, recovering and generalising \cref{thm:10:kedlaya-liu} --- and up to the
final Fredholm step, no analysis was done at all.
\end{remark}
